% ─── Chapter 1: Introduction ───────────────────────────────────────
\chapter{Introduction}

\section{Problem Statement}

The goal of this project is to develop an AI agent capable of completing a
\emph{Pok\'{e}mon Nuzlocke challenge} using deep reinforcement learning.
The chosen game is \emph{Pok\'{e}mon Brilliant Diamond} (BDSP), and the agent
must defeat all mandatory trainers without losing---subject to constraints that
transform standard battles into a complex sequential decision-making problem
under uncertainty with irreversible consequences.

A Nuzlocke challenge is a self-imposed difficulty rule set widely adopted by the
Pok\'{e}mon community.  Its core constraints are:

\begin{itemize}
  \item \textbf{First-catch rule:} Only the first Pok\'{e}mon encountered in
        each new location may be captured.
  \item \textbf{Permadeath:} If a Pok\'{e}mon faints during battle, it is
        considered permanently lost and cannot be used again.
  \item \textbf{Run failure:} If the player loses a battle with no usable
        Pok\'{e}mon remaining, the entire run ends and must restart.
\end{itemize}

These rules introduce permanent consequences and limited resources, turning a
turn-based strategy game into a problem that requires long-horizon planning,
risk assessment, and resource management---all qualities that make it a
compelling testbed for reinforcement learning.

\section{Stakeholder Context}

This project is carried out as part of the \textbf{DSPRO2} course at the
Hochschule Luzern (HSLU) during the Spring 2026 semester.  Beyond fulfilling
academic requirements, the project serves as a \emph{reusable experimental
benchmark} for evaluating reinforcement learning algorithms on generalisation,
robustness, and risk-aware decision-making.

The primary stakeholders are:

\begin{itemize}
  \item \textbf{Course supervisors:} Expect a professional-grade report
        demonstrating an end-to-end ML workflow, experiment tracking, honest
        validation, and actionable conclusions.
  \item \textbf{Project team:} Uses the system as a portfolio-quality technical
        artifact proving the ability to move from prototype to a
        production-ready mindset.
  \item \textbf{RL research community:} The Nuzlocke benchmark, if
        successful, provides a challenging, well-defined testbed for
        comparing RL approaches on tasks with sparse rewards, partial
        observability, and irreversible actions.
\end{itemize}

\section{Objectives \& Success Metrics}

The overarching objective is to train an AI agent that successfully completes a
Nuzlocke gauntlet---a predefined sequence of battles against trainer teams from
Pok\'{e}mon Brilliant Diamond.

Success is measured by the following metrics:

\begin{center}
\small
\begin{tabularx}{\textwidth}{l X}
  \toprule
  \textbf{Metric} & \textbf{Description} \\
  \midrule
  Gauntlet completion rate  & Fraction of runs where the agent beats all
                             trainers in the correct sequence. \\
  Progression depth         & Average number of battles completed before
                             failure. \\
  Remaining resources       & HP and Pok\'{e}mon still available upon
                             gauntlet completion. \\
  Generalisation            & Performance on unseen gauntlet orderings or
                             opponent teams without further training. \\
  Model efficiency          & Minimum model size (parameters) needed to reach
                             a given competency level. \\
  \bottomrule
\end{tabularx}
\end{center}

Near-term success is defined as achieving a stable win rate above 50\,\%
against heuristic opponents in the medium curriculum stage (see
Chapter~\ref{ch:rl_framework}).

\section{Project Scope}

The project focuses exclusively on \textbf{battle decision-making}:

\begin{itemize}
  \item Move selection and Pok\'{e}mon switching during battles.
  \item Team composition (in later milestones).
  \item Sequential gauntlet management with Nuzlocke constraints.
\end{itemize}

The following are \textbf{explicitly excluded}:

\begin{itemize}
  \item Game-world exploration and traversal.
  \item Interaction with an actual game client---the system operates entirely
        through the Pok\'{e}mon Showdown battle simulator.
  \item Multiplayer or competitive battling against human players.
\end{itemize}

\section{Document Structure}

This report is organised as follows.  Chapter~\ref{ch:game_env} provides
background on Pok\'{e}mon battle mechanics and the simulation environment.
Chapter~\ref{ch:challenge} formalises the challenge model, data provenance,
and milestone roadmap.  Chapter~\ref{ch:rl_framework} describes the
reinforcement learning formulation, PPO algorithm, reward design, and
curriculum learning setup.  Chapter~\ref{ch:architecture} details the system
architecture---observation pipeline, model, action space, and self-play
mechanism.  Chapter~\ref{ch:infrastructure} covers the distributed training
infrastructure and experiment tracking.  Chapter~\ref{ch:experiments} presents
the iterative development process through multiple experimental phases.
Chapter~\ref{ch:results} analyses training results.  Chapter~\ref{ch:risk}
discusses risk, ethical, and human-in-the-loop considerations.
Chapter~\ref{ch:conclusions} concludes with lessons learned and next steps.
Chapter~\ref{ch:reproduction} provides a complete reproduction guide.
